\section{讨论、局限性与可复现性}
\label{sec:discussion}

\subsection{经验教训}

这次“设计良好却在实测中退化”的经历，给出几条超出本系统的教训。

\textbf{教训一：分层决策必须内建向上移交的强制条件。} 分层的初衷是“廉价层处理简单情形、昂贵层处理复杂情形”，但若廉价层对高频情形总能给出一个自信的答案，分层就退化为“廉价层独裁”。有效的分层必须让廉价层\emph{能够并且被迫承认不确定}——例如以“连续 $k$ 次动作无进展”作为强制回退到推理层的条件。缺了这一环，越廉价、越自信的规则层，越会把智能层饿死。

\textbf{教训二：动作必须有有效性反馈，而非仅有状态触发。} 现有系统是纯粹的“状态$\to$动作”反射：见到空闲提示符就发“继续”。它缺少“动作$\to$效果”的闭环——没有观察“继续”是否真的带来了进展。任何自动化系统一旦缺失效果反馈，就无法区分“有用的重复”与“空转的重复”，从而稳定地滑向后者。

\textbf{教训三：评估自主系统要以有效性指标为主、效率指标为辅。} 如 \S\ref{sec:eval} 所示，效率指标（命中率、延迟、成本、崩溃率）会在系统空转时集体“亮绿灯”。只有度量“任务是否真正推进”的有效性指标，才能戳破“高效空转”的假象。

\textbf{教训四：不报错的失败比报错的失败更危险。} 系统三天里没有崩溃、没有异常，却也没做任何有用的事。这种“沉默的失败”难以通过监控告警发现，需要专门的进展度量与人工抽检才能暴露。

\subsection{设计权衡与修复}

\S\ref{sec:rootcause} 已指出修复方向：为“继续”类动作引入基于屏幕实质变化或 Hook 工具调用的有效性反馈；设定规则层向 LLM 层的强制移交条件；区分“语法空闲”与“语义空闲”。这些修复的共同代价是引入更多状态跟踪与可能的 LLM 调用，从而抬高成本——但这正是恢复“智能”所必须支付的价格。当前实现选择了成本最优（近零 LLM），代价是有效性归零；合理的设计应在二者间取得平衡，而非把 $\rho$ 最大化本身当作目标。

\subsection{安全边界与失败模式}

自动化操作终端本质具有风险：系统代替人点击“确认”“继续”，误判即误操作。系统设置了 Token 鉴权、本机回环、会话名清洗、关键阶段强制人工确认、自主执行隔离分支等边界。本次部署未观察到误操作\emph{恰恰是因为它几乎没做实质操作}——这提示我们，一旦有效性提升、系统真正开始自主行动，安全边界将面临远比现在更严峻的考验。

\subsection{局限性}

其一，本文的实证基于单次、单机、46 次决策的部署，样本有限，结论定位为“有根因支撑的负面案例”而非普适量化率。其二，规则状态机与所用 CLI 的界面耦合，CLI 升级可能改变退化的具体形态。其三，系统的自主执行引擎核心以\emph{闭源}形式分发，开源包不含该核心而降级运行，这对完全基于开源代码的复现构成约束。

\subsection{可复现性说明}

前端、会话管理、感知层、决策层（含本文剖析的规则状态机）与领域策略框架均在开源仓库中可查证。本文的关键证据——46 次决策的历史记录与引擎输出文件——由系统运行时自动持久化于 SQLite 与临时文件，分析脚本仅做只读统计，可在相同数据上复现。三层自主循环的编排逻辑开源可见，但其执行引擎核心闭源，故涉及“自主交付是否成功”的复现需依赖正式发行版。为保证可复现，应记录仓库地址、版本 tag、CLI 与模型版本及运行平台。
